

\documentclass[11pt,respuestas,a4]{aleph-examen}
%\documentclass[10pt,a4]{aleph-examen}

% -- Paquetes adicionales 
\usepackage{aleph-comandos}
\usepackage{booktabs}
\usepackage{multicol}
\usepackage{enumitem}
\usepackage{array}
\usepackage{amsmath}

% -- Datos 
\institucion{Facultad de Ciencias Exactas, Naturales y Ambientales}
\carrera{Catálogo STEM}
\asignatura{Métodos Numéricos}
\tema{Examen}
\autor{Mario Cueva - Andrés Merino}
\fecha{Periodo 2026-1}

\logouno[0.18\textwidth]{Logos/logoPUCE_04_ac}
\definecolor{colortext}{HTML}{1957d1}
\definecolor{colordef}{HTML}{1957d1}
\fuente{montserrat}

\begin{document}

\encabezado

\section*{Indicaciones}
\begin{itemize}
    \item El examen tiene una duración de \textbf{90 minutos}.
    \item Se permite únicamente el uso de \textbf{calculadora científica}, lápiz y papel.
    \item Justifique todos los procedimientos. Las respuestas sin desarrollo tendrán puntaje parcial limitado.
    \item Redondee los resultados finales a \textbf{cuatro cifras decimales}, salvo que se indique lo contrario.
    \item En los métodos iterativos, presente las iteraciones en tablas ordenadas.
    \item En los ejercicios de integración y diferenciación, complete las tablas cuando sea necesario.
\end{itemize}

\section*{Examen}

\begin{preguntas}

%%%%%%%%%%%%%%%%%%%%%%%%%%%%%%%%%%%%%%%%
% Pregunta 1
%%%%%%%%%%%%%%%%%%%%%%%%%%%%%%%%%%%%%%%%

\item
Considere el siguiente sistema lineal: \puntaje{8}
\[
\begin{cases}
6x+y=14,\\
x+5y=13.
\end{cases}
\]

Aplique \textbf{tres iteraciones} de los métodos de \textbf{Jacobi} y \textbf{Gauss-Seidel}, partiendo de:
\[
(x^{(0)},y^{(0)})=(0,0).
\]

\begin{enumerate}[label=\alph*)]
    \item Despeje las ecuaciones para aplicar ambos métodos.
    \item Complete la tabla para el método de Jacobi.
    \item Complete la tabla para el método de Gauss-Seidel.
    \item Compare brevemente las aproximaciones obtenidas.
\end{enumerate}

\[
x=\underline{\hspace{4cm}}, 
\qquad 
y=\underline{\hspace{4cm}}
\]

\begin{center}
\begin{minipage}{0.47\textwidth}
\centering
\textbf{Método de Jacobi}

\renewcommand{\arraystretch}{1.3}
\begin{tabular}{c|c|c}
\toprule
$k$ & $x^{(k)}$ & $y^{(k)}$ \\
\midrule
0 & 0.0000 & 0.0000 \\
1 &  &  \\
2 &  &  \\
3 &  &  \\
\bottomrule
\end{tabular}
\end{minipage}
\hfill
\begin{minipage}{0.47\textwidth}
\centering
\textbf{Método de Gauss-Seidel}

\renewcommand{\arraystretch}{1.3}
\begin{tabular}{c|c|c}
\toprule
$k$ & $x^{(k)}$ & $y^{(k)}$ \\
\midrule
0 & 0.0000 & 0.0000 \\
1 &  &  \\
2 &  &  \\
3 &  &  \\
\bottomrule
\end{tabular}
\end{minipage}
\end{center}

Comparación:

\[
\underline{\hspace{14cm}}
\]

\[
\underline{\hspace{14cm}}
\]

\begin{respuesta}
Primero despejamos cada variable:
\[
x=\frac{14-y}{6},
\qquad
y=\frac{13-x}{5}.
\]

\textbf{Método de Jacobi.} En Jacobi se usan los valores de la iteración anterior:
\[
x^{(k+1)}=\frac{14-y^{(k)}}{6},
\qquad
y^{(k+1)}=\frac{13-x^{(k)}}{5}.
\]

\begin{center}
\renewcommand{\arraystretch}{1.3}
\begin{tabular}{c|c|c}
\toprule
$k$ & $x^{(k)}$ & $y^{(k)}$ \\
\midrule
0 & 0.0000 & 0.0000 \\
1 & 2.3333 & 2.6000 \\
2 & 1.9000 & 2.1333 \\
3 & 1.9778 & 2.2200 \\
\bottomrule
\end{tabular}
\end{center}

\textbf{Método de Gauss-Seidel.} En Gauss-Seidel se usan inmediatamente los valores actualizados:
\[
x^{(k+1)}=\frac{14-y^{(k)}}{6},
\qquad
y^{(k+1)}=\frac{13-x^{(k+1)}}{5}.
\]

\begin{center}
\renewcommand{\arraystretch}{1.3}
\begin{tabular}{c|c|c}
\toprule
$k$ & $x^{(k)}$ & $y^{(k)}$ \\
\midrule
0 & 0.0000 & 0.0000 \\
1 & 2.3333 & 2.1333 \\
2 & 1.9778 & 2.2044 \\
3 & 1.9659 & 2.2068 \\
\bottomrule
\end{tabular}
\end{center}

Por tanto, después de tres iteraciones:
\[
\boxed{\text{Jacobi: }(x,y)\approx(1.9778,2.2200)}
\]
\[
\boxed{\text{Gauss-Seidel: }(x,y)\approx(1.9659,2.2068)}
\]

En este caso, Gauss-Seidel se acerca más rápido a la solución porque actualiza cada variable dentro de la misma iteración.
\end{respuesta}

%%%%%%%%%%%%%%%%%%%%%%%%%%%%%%%%%%%%%%%%
% Pregunta 2
%%%%%%%%%%%%%%%%%%%%%%%%%%%%%%%%%%%%%%%%

\item
Construya el polinomio interpolante cuadrático de Newton para los puntos: \puntaje{6}
\[
(0,2),\qquad (1,3),\qquad (3,23).
\]

Luego estime el valor de la función en:
\[
x=2.
\]

Complete la tabla de diferencias divididas:

\begin{center}
\renewcommand{\arraystretch}{1.5}
\begin{tabular}{c|c|c|c}
\toprule
$i$ & $x_i$ & $f(x_i)$ & Diferencias divididas \\
\midrule
0 & 0 & 2 &  \\
1 & 1 & 3 &  \\
2 & 3 & 23 &  \\
\bottomrule
\end{tabular}
\end{center}

Escriba el polinomio de Newton:

\[
P_2(x)=\underline{\hspace{10cm}}
\]

Evalúe:

\[
P_2(2)=\underline{\hspace{5cm}}
\]

\begin{respuesta}
Sean:
\[
x_0=0,\quad f_0=2,
\qquad
x_1=1,\quad f_1=3,
\qquad
x_2=3,\quad f_2=23.
\]

Las diferencias divididas de primer orden son:
\[
f[x_0,x_1]=\frac{3-2}{1-0}=1,
\]
\[
f[x_1,x_2]=\frac{23-3}{3-1}=10.
\]

La diferencia dividida de segundo orden es:
\[
f[x_0,x_1,x_2]=\frac{10-1}{3-0}=3.
\]

Por tanto, el polinomio de Newton es:
\[
P_2(x)=2+1(x-0)+3(x-0)(x-1).
\]

Simplificando:
\[
P_2(x)=2+x+3x(x-1).
\]

\[
P_2(x)=2+x+3x^2-3x.
\]

\[
P_2(x)=3x^2-2x+2.
\]

Ahora evaluamos en $x=2$:
\[
P_2(2)=3(2)^2-2(2)+2.
\]

\[
P_2(2)=12-4+2=10.
\]

Por lo tanto:
\[
\boxed{P_2(x)=3x^2-2x+2}
\]
\[
\boxed{P_2(2)=10.0000}
\]
\end{respuesta}

%%%%%%%%%%%%%%%%%%%%%%%%%%%%%%%%%%%%%%%%
% Pregunta 3
%%%%%%%%%%%%%%%%%%%%%%%%%%%%%%%%%%%%%%%%

\item
Ajuste una recta de regresión lineal de la forma: \puntaje{6}
\[
y=a_0+a_1x
\]
para los siguientes datos:

\[
\begin{array}{c|ccccc}
x_i & 1 & 2 & 3 & 4 & 5 \\
\hline
y_i & 2 & 5 & 7 & 10 & 11
\end{array}
\]

\begin{enumerate}[label=\alph*)]
    \item Complete la tabla de cálculo.
    \item Determine los coeficientes $a_0$ y $a_1$.
    \item Escriba el modelo ajustado.
    \item Estime el valor de $y$ cuando $x=6$.
\end{enumerate}

\begin{center}
\renewcommand{\arraystretch}{1.4}
\begin{tabular}{c|c|c|c}
\toprule
$x_i$ & $y_i$ & $x_i^2$ & $x_iy_i$ \\
\midrule
1 & 2 &  &  \\
2 & 5 &  &  \\
3 & 7 &  &  \\
4 & 10 &  &  \\
5 & 11 &  &  \\
\midrule
$\sum$ &  &  &  \\
\bottomrule
\end{tabular}
\end{center}

\[
a_1=\underline{\hspace{5cm}}, 
\qquad
a_0=\underline{\hspace{5cm}}
\]

Modelo:
\[
y=\underline{\hspace{7cm}}
\]

Estimación:
\[
y(6)=\underline{\hspace{5cm}}
\]

\begin{respuesta}
Construimos la tabla de cálculo:

\begin{center}
\renewcommand{\arraystretch}{1.4}
\begin{tabular}{c|c|c|c}
\toprule
$x_i$ & $y_i$ & $x_i^2$ & $x_iy_i$ \\
\midrule
1 & 2 & 1 & 2 \\
2 & 5 & 4 & 10 \\
3 & 7 & 9 & 21 \\
4 & 10 & 16 & 40 \\
5 & 11 & 25 & 55 \\
\midrule
15 & 35 & 55 & 128 \\
\bottomrule
\end{tabular}
\end{center}

Por tanto:
\[
n=5,
\qquad
\sum x_i=15,
\qquad
\sum y_i=35,
\]
\[
\sum x_i^2=55,
\qquad
\sum x_iy_i=128.
\]

La pendiente se calcula con:
\[
a_1=\frac{n\sum x_iy_i-\sum x_i\sum y_i}{n\sum x_i^2-(\sum x_i)^2}.
\]

Sustituyendo:
\[
a_1=\frac{5(128)-15(35)}{5(55)-15^2}.
\]

\[
a_1=\frac{640-525}{275-225}
=\frac{115}{50}=2.3.
\]

Además:
\[
\bar{x}=\frac{15}{5}=3,
\qquad
\bar{y}=\frac{35}{5}=7.
\]

Entonces:
\[
a_0=\bar{y}-a_1\bar{x}.
\]

\[
a_0=7-2.3(3)=7-6.9=0.1.
\]

Por tanto, el modelo ajustado es:
\[
\boxed{y=0.1+2.3x}.
\]

Para $x=6$:
\[
y=0.1+2.3(6)=0.1+13.8=13.9.
\]

Por lo tanto:
\[
\boxed{y(6)\approx 13.9000}.
\]
\end{respuesta}

%%%%%%%%%%%%%%%%%%%%%%%%%%%%%%%%%%%%%%%%
% Pregunta 4
%%%%%%%%%%%%%%%%%%%%%%%%%%%%%%%%%%%%%%%%

\item
Para los datos: \puntaje{6}
\[
\begin{array}{c|cccc}
x_i & 0 & 1 & 2 & 3 \\
\hline
y_i & 2 & 4 & 9 & 17
\end{array}
\]

ajuste un modelo cuadrático de la forma:
\[
y=a_0+a_1x+a_2x^2
\]
utilizando la forma matricial de mínimos cuadrados.

\begin{enumerate}[label=\alph*)]
    \item Escriba la matriz de diseño $A$.
    \item Escriba el vector $\mathbf{y}$.
    \item Plantee el sistema normal $A^TA\mathbf{a}=A^T\mathbf{y}$.
    \item Determine los coeficientes del modelo.
\end{enumerate}

\[
A=
\begin{pmatrix}
\hspace{2.2cm} . & \hspace{1.2cm} & \hspace{1.2cm} \\[0.4cm]
\hspace{1.2cm} & \hspace{1.2cm} & \hspace{1.2cm} \\[0.4cm]
\hspace{1.2cm} & \hspace{1.2cm} & \hspace{1.2cm} \\[0.4cm]
\hspace{1.2cm} & \hspace{1.2cm} & \hspace{1.2cm}
\end{pmatrix},
\qquad
\mathbf{y}=
\begin{pmatrix}
\hspace{1.2cm} .\\[0.4cm]
\hspace{1.2cm}\\[0.4cm]
\hspace{1.2cm}\\[0.4cm]
\hspace{1.2cm}
\end{pmatrix}
\]

\[
A^TA=
\begin{pmatrix}
\hspace{2.2cm} .& \hspace{1.2cm} & \hspace{1.2cm} \\[0.4cm]
\hspace{1.2cm} & \hspace{1.2cm} & \hspace{1.2cm} \\[0.4cm]
\hspace{1.2cm} & \hspace{1.2cm} & \hspace{1.2cm}
\end{pmatrix},
\qquad
A^T\mathbf{y}=
\begin{pmatrix}
\hspace{1.2cm}.\\[0.4cm]
\hspace{1.2cm}\\[0.4cm]
\hspace{1.2cm}
\end{pmatrix}
\]

Sistema normal:

\[
\begin{pmatrix}
\hspace{2.2cm} .& \hspace{1.2cm} & \hspace{1.2cm} \\[0.4cm]
\hspace{1.2cm} & \hspace{1.2cm} & \hspace{1.2cm} \\[0.4cm]
\hspace{1.2cm} & \hspace{1.2cm} & \hspace{1.2cm}
\end{pmatrix}
\begin{pmatrix}
a_0\\
a_1\\
a_2
\end{pmatrix}
=
\begin{pmatrix}
\hspace{1.2cm}.\\[0.4cm]
\hspace{1.2cm}\\[0.4cm]
\hspace{1.2cm}
\end{pmatrix}
\]

Modelo ajustado:

\[
y=\underline{\hspace{8cm}}
\]

\begin{respuesta}
La matriz de diseño para el modelo cuadrático es:
\[
A=
\begin{pmatrix}
1 & 0 & 0 \\
1 & 1 & 1 \\
1 & 2 & 4 \\
1 & 3 & 9
\end{pmatrix},
\qquad
\mathbf{y}=
\begin{pmatrix}
2\\
4\\
9\\
17
\end{pmatrix}.
\]

El sistema normal es:
\[
A^TA\mathbf{a}=A^T\mathbf{y}.
\]

Calculamos:
\[
A^TA=
\begin{pmatrix}
4 & 6 & 14 \\
6 & 14 & 36 \\
14 & 36 & 98
\end{pmatrix}.
\]

Además:
\[
A^T\mathbf{y}=
\begin{pmatrix}
2+4+9+17\\
0(2)+1(4)+2(9)+3(17)\\
0^2(2)+1^2(4)+2^2(9)+3^2(17)
\end{pmatrix}.
\]

\[
A^T\mathbf{y}=
\begin{pmatrix}
32\\
73\\
193
\end{pmatrix}.
\]

Por tanto, el sistema normal queda:
\[
\begin{pmatrix}
4 & 6 & 14 \\
6 & 14 & 36 \\
14 & 36 & 98
\end{pmatrix}
\begin{pmatrix}
a_0\\
a_1\\
a_2
\end{pmatrix}
=
\begin{pmatrix}
32\\
73\\
193
\end{pmatrix}.
\]

Resolviendo el sistema se obtiene:
\[
a_0=2,
\qquad
a_1=0.5,
\qquad
a_2=1.5.
\]

Así, el modelo cuadrático ajustado es:
\[
\boxed{y=2+0.5x+1.5x^2}.
\]
\end{respuesta}

%%%%%%%%%%%%%%%%%%%%%%%%%%%%%%%%%%%%%%%%
% Pregunta 5
%%%%%%%%%%%%%%%%%%%%%%%%%%%%%%%%%%%%%%%%

\item
Utilice diferenciación numérica para aproximar $f'(1)$ de la función: \puntaje{6}
\[
f(x)=\ln(x+2)
\]
con $h=0.2$, aplicando:

\begin{enumerate}[label=\alph*)]
    \item Diferencia hacia adelante.
    \item Diferencia hacia atrás.
    \item Diferencia centrada.
    \item Compare brevemente los resultados con el valor exacto:
    \[
    f'(1)=\frac{1}{3}.
    \]
\end{enumerate}

\begin{center}
\begin{minipage}{0.45\textwidth}
\centering
\renewcommand{\arraystretch}{1.4}
\begin{tabular}{c|c}
\toprule
$x$ & $f(x)=\ln(x+2)$ \\
\midrule
0.8 &  \\
1.0 &  \\
1.2 &  \\
\bottomrule
\end{tabular}
\end{minipage}
\hfill
\begin{minipage}{0.48\textwidth}
\[
f'(1)_{\text{adelante}}\approx \underline{\hspace{3cm}}
\]
\[
f'(1)_{\text{atrás}}\approx \underline{\hspace{3cm}}
\]
\[
f'(1)_{\text{centrada}}\approx \underline{\hspace{3cm}}
\]
\end{minipage}
\end{center}

Comparación:

\[
\underline{\hspace{14cm}}
\]

\[
\underline{\hspace{14cm}}
\]

\begin{respuesta}
Se tiene:
\[
f(x)=\ln(x+2),
\qquad
x=1,
\qquad
h=0.2.
\]

Calculamos los valores necesarios:
\[
f(0.8)=\ln(2.8)=1.0296,
\]
\[
f(1)=\ln(3)=1.0986,
\]
\[
f(1.2)=\ln(3.2)=1.1632.
\]

\begin{center}
\renewcommand{\arraystretch}{1.4}
\begin{tabular}{c|c}
\toprule
$x$ & $f(x)=\ln(x+2)$ \\
\midrule
0.8 & 1.0296 \\
1.0 & 1.0986 \\
1.2 & 1.1632 \\
\bottomrule
\end{tabular}
\end{center}

\textbf{a) Diferencia hacia adelante:}
\[
f'(1)\approx \frac{f(1.2)-f(1)}{0.2}.
\]

\[
f'(1)\approx \frac{1.1632-1.0986}{0.2}
=\frac{0.0646}{0.2}=0.3230.
\]

\textbf{b) Diferencia hacia atrás:}
\[
f'(1)\approx \frac{f(1)-f(0.8)}{0.2}.
\]

\[
f'(1)\approx \frac{1.0986-1.0296}{0.2}
=\frac{0.0690}{0.2}=0.3450.
\]

\textbf{c) Diferencia centrada:}
\[
f'(1)\approx \frac{f(1.2)-f(0.8)}{2(0.2)}.
\]

\[
f'(1)\approx \frac{1.1632-1.0296}{0.4}
=\frac{0.1336}{0.4}=0.3340.
\]

El valor exacto es:
\[
f'(1)=\frac{1}{3}=0.3333.
\]

Por lo tanto:
\[
\boxed{\text{Adelante: }0.3230}
\]
\[
\boxed{\text{Atrás: }0.3450}
\]
\[
\boxed{\text{Centrada: }0.3340}
\]

La diferencia centrada es la aproximación más cercana al valor exacto.
\end{respuesta}

%%%%%%%%%%%%%%%%%%%%%%%%%%%%%%%%%%%%%%%%
% Pregunta 6
%%%%%%%%%%%%%%%%%%%%%%%%%%%%%%%%%%%%%%%%

\item
Aproxime la integral: \puntaje{6}
\[
\int_0^2 (x^3+1)\,dx
\]
mediante la \textbf{regla del trapecio compuesta} con $n=4$ subintervalos, luego calcule el error absoluto.

\begin{enumerate}[label=\alph*)]
    \item Calcule el valor de $h$.
    \item Complete la tabla de valores.
    \item Aplique la fórmula del trapecio compuesto.
    \item Calcule el error absoluto.
\end{enumerate}

\[
h=\underline{\hspace{5cm}}
\]

\begin{center}
\renewcommand{\arraystretch}{1.4}
\begin{tabular}{c|c}
\toprule
$x_i$ & $f(x_i)=x_i^3+1$ \\
\midrule
0.0 &  \\
0.5 &  \\
1.0 &  \\
1.5 &  \\
2.0 &  \\
\bottomrule
\end{tabular}
\end{center}

Aplicación de la fórmula:

\[
\int_0^2 (x^3+1)\,dx
\approx
\underline{\hspace{8cm}}
\]

Error absoluto:

\[
E_a=\underline{\hspace{6cm}}
\]

\begin{respuesta}
Primero calculamos:
\[
h=\frac{b-a}{n}=\frac{2-0}{4}=0.5.
\]

Los nodos son:
\[
x_0=0,
\quad
x_1=0.5,
\quad
x_2=1,
\quad
x_3=1.5,
\quad
x_4=2.
\]

Evaluamos $f(x)=x^3+1$:
\[
\begin{array}{c|ccccc}
x_i & 0 & 0.5 & 1 & 1.5 & 2 \\
\hline
f(x_i) & 1 & 1.125 & 2 & 4.375 & 9
\end{array}
\]

La regla del trapecio compuesta es:
\[
\int_a^b f(x)\,dx
\approx
\frac{h}{2}
\left[
f(x_0)+2\sum_{i=1}^{n-1}f(x_i)+f(x_n)
\right].
\]

Sustituyendo:
\[
\int_0^2 (x^3+1)\,dx
\approx
\frac{0.5}{2}
\left[
1+2(1.125+2+4.375)+9
\right].
\]

\[
=
0.25
\left[
1+2(7.5)+9
\right].
\]

\[
=
0.25(25)=6.25.
\]

El valor exacto es:
\[
6.
\]

El error absoluto es:
\[
E_a=|6-6.25|=0.25.
\]

Por lo tanto:
\[
\boxed{\int_0^2 (x^3+1)\,dx\approx 6.2500}
\]
\[
\boxed{E_a=0.2500}.
\]
\end{respuesta}

%%%%%%%%%%%%%%%%%%%%%%%%%%%%%%%%%%%%%%%%
% Pregunta 7
%%%%%%%%%%%%%%%%%%%%%%%%%%%%%%%%%%%%%%%%

\item
Aproxime la integral: \puntaje{6}
\[
\int_0^2 (x^3-2x+4)\,dx
\]
mediante la \textbf{regla de Simpson $1/3$ simple}.

\begin{enumerate}[label=\alph*)]
    \item Identifique los valores de $a$, $b$ y el punto medio $m$.
    \item Evalúe la función en los tres puntos necesarios.
    \item Aplique la fórmula de Simpson $1/3$ simple.
    \item Indique si el resultado puede ser exacto para este tipo de función y explique brevemente.
\end{enumerate}

\[
a=\underline{\hspace{3cm}},
\qquad
b=\underline{\hspace{3cm}},
\qquad
m=\underline{\hspace{3cm}}
\]

\begin{center}
\renewcommand{\arraystretch}{1.4}
\begin{tabular}{c|c}
\toprule
$x$ & $f(x)=x^3-2x+4$ \\
\midrule
$a=$ &  \\
$m=$ &  \\
$b=$ &  \\
\bottomrule
\end{tabular}
\end{center}

Resultado:

\[
\int_0^2 (x^3-2x+4)\,dx
\approx
\underline{\hspace{7cm}}
\]

Explicación:

\[
\underline{\hspace{14cm}}
\]
\[
\underline{\hspace{14cm}}
\]

\begin{respuesta}
Para Simpson $1/3$ simple se usa:
\[
\int_a^b f(x)\,dx
\approx
\frac{b-a}{6}
\left[
f(a)+4f(m)+f(b)
\right],
\]
donde:
\[
m=\frac{a+b}{2}.
\]

En este caso:
\[
a=0,
\qquad
b=2,
\qquad
m=\frac{0+2}{2}=1.
\]

La función es:
\[
f(x)=x^3-2x+4.
\]

Calculamos:
\[
f(0)=0^3-2(0)+4=4,
\]
\[
f(1)=1^3-2(1)+4=3,
\]
\[
f(2)=2^3-2(2)+4=8.
\]

Aplicando Simpson $1/3$ simple:
\[
\int_0^2 (x^3-2x+4)\,dx
\approx
\frac{2-0}{6}
\left[
4+4(3)+8
\right].
\]

\[
=
\frac{2}{6}
\left[
4+12+8
\right].
\]

\[
=
\frac{2}{6}(24)=8.
\]

Por tanto:
\[
\boxed{\int_0^2 (x^3-2x+4)\,dx\approx 8.0000}.
\]

En este caso, Simpson $1/3$ simple da el valor exacto porque la función es un polinomio de grado $3$.
\end{respuesta}

%%%%%%%%%%%%%%%%%%%%%%%%%%%%%%%%%%%%%%%%
% Pregunta 8
%%%%%%%%%%%%%%%%%%%%%%%%%%%%%%%%%%%%%%%%

\item
Aproxime la integral: \puntaje{8}
\[
\int_0^{\frac{\pi}{2}} \sin(x)\,dx
\]
mediante la \textbf{regla de Simpson $1/3$ compuesta} con $n=6$ subintervalos.

Además, calcule el error absoluto, considerando que el valor exacto de la integral es:
\[
1.
\]

\begin{enumerate}[label=\alph*)]
    \item Calcule el valor de $h$.
    \item Complete la tabla de valores.
    \item Identifique los valores que se multiplican por $4$ y por $2$.
    \item Aplique la fórmula de Simpson $1/3$ compuesta.
    \item Calcule el error absoluto.
\end{enumerate}

\[
h=\underline{\hspace{5cm}}
\]

\begin{center}
\renewcommand{\arraystretch}{1.4}
\begin{tabular}{c|c|c}
\toprule
$i$ & $x_i$ & $f(x_i)=\sin(x_i)$ \\
\midrule
0 & $0$ &  \\
1 & $\frac{\pi}{12}$ &  \\
2 & $\frac{\pi}{6}$ &  \\
3 & $\frac{\pi}{4}$ &  \\
4 & $\frac{\pi}{3}$ &  \\
5 & $\frac{5\pi}{12}$ &  \\
6 & $\frac{\pi}{2}$ &  \\
\bottomrule
\end{tabular}
\end{center}

Valores con coeficiente $4$: $
\underline{\hspace{10cm}}
$

\vspace{1cm}
Valores con coeficiente $2$: $
\underline{\hspace{10cm}}
$

Resultado:

\[
\int_0^{\frac{\pi}{2}} \sin(x)\,dx
\approx
\underline{\hspace{7cm}}
\]

Error absoluto: $
E_a=\underline{\hspace{6cm}}
$

\begin{respuesta}
Primero calculamos:
\[
h=\frac{b-a}{n}
=
\frac{\frac{\pi}{2}-0}{6}
=
\frac{\pi}{12}.
\]

En forma decimal:
\[
h\approx 0.2618.
\]

Los nodos son:
\[
x_0=0,
\quad
x_1=\frac{\pi}{12},
\quad
x_2=\frac{\pi}{6},
\quad
x_3=\frac{\pi}{4},
\quad
x_4=\frac{\pi}{3},
\quad
x_5=\frac{5\pi}{12},
\quad
x_6=\frac{\pi}{2}.
\]

Evaluamos $f(x)=\sin(x)$:

\[
\begin{array}{c|c|c}
i & x_i & f(x_i)=\sin(x_i) \\
\hline
0 & 0 & 0.0000 \\
1 & \frac{\pi}{12} & 0.2588 \\
2 & \frac{\pi}{6} & 0.5000 \\
3 & \frac{\pi}{4} & 0.7071 \\
4 & \frac{\pi}{3} & 0.8660 \\
5 & \frac{5\pi}{12} & 0.9659 \\
6 & \frac{\pi}{2} & 1.0000
\end{array}
\]

La regla de Simpson $1/3$ compuesta es:
\[
\int_a^b f(x)\,dx
\approx
\frac{h}{3}
\left[
f(x_0)
+
4\sum_{\substack{i=1\\i\text{ impar}}}^{n-1}f(x_i)
+
2\sum_{\substack{i=2\\i\text{ par}}}^{n-2}f(x_i)
+
f(x_n)
\right].
\]

En este caso, los índices impares son:
\[
i=1,3,5.
\]

Por tanto, los valores que se multiplican por $4$ son:
\[
f(x_1)=0.2588,
\qquad
f(x_3)=0.7071,
\qquad
f(x_5)=0.9659.
\]

Los índices pares interiores son:
\[
i=2,4.
\]

Por tanto, los valores que se multiplican por $2$ son:
\[
f(x_2)=0.5000,
\qquad
f(x_4)=0.8660.
\]

Sustituyendo:
\[
\int_0^{\frac{\pi}{2}}\sin(x)\,dx
\approx
\frac{0.2618}{3}
\left[
0.0000
+
4(0.2588+0.7071+0.9659)
+
2(0.5000+0.8660)
+
1.0000
\right].
\]

\[
\int_0^{\frac{\pi}{2}}\sin(x)\,dx
\approx
0.0873
\left[
4(1.9318)+2(1.3660)+1.0000
\right].
\]

\[
\int_0^{\frac{\pi}{2}}\sin(x)\,dx
\approx
0.0873
\left[
7.7272+2.7320+1.0000
\right].
\]

\[
\int_0^{\frac{\pi}{2}}\sin(x)\,dx
\approx
0.0873(11.4592).
\]

\[
\int_0^{\frac{\pi}{2}}\sin(x)\,dx
\approx
1.0000.
\]

El valor exacto es:
\[
1.
\]

Entonces, el error absoluto es:
\[
E_a=|1-1.0000|=0.0000.
\]

Por tanto:
\[
\boxed{\int_0^{\frac{\pi}{2}}\sin(x)\,dx\approx 1.0000}
\]
\[
\boxed{E_a\approx 0.0000}.
\]

El resultado es muy cercano al valor exacto porque Simpson $1/3$ compuesta suele dar buenas aproximaciones cuando la función es suave.
\end{respuesta}

%%%%%%%%%%%%%%%%%%%%%%%%%%%%%%%%%%%%%%%%
% Pregunta 9
%%%%%%%%%%%%%%%%%%%%%%%%%%%%%%%%%%%%%%%%

\item
Pregunta conceptual de integración numérica: \puntaje{3}

Explique una diferencia entre la regla del trapecio compuesta y la regla de Simpson $1/3$ compuesta. Además, indique cuál de las dos suele dar una mejor aproximación cuando la función es suave y se usa el mismo número de subintervalos.

\[
\underline{\hspace{14cm}}
\]

\[
\underline{\hspace{14cm}}
\]

\[
\underline{\hspace{14cm}}
\]

\begin{respuesta}
La regla del trapecio compuesta aproxima el área bajo la curva usando varios trapecios, es decir, aproxima la función con segmentos rectos en cada subintervalo.

En cambio, la regla de Simpson $1/3$ compuesta aproxima la función mediante parábolas, agrupando los subintervalos de dos en dos.

Cuando la función es suave y se usa un número adecuado de subintervalos, la regla de Simpson $1/3$ compuesta suele dar una mejor aproximación que la regla del trapecio compuesta, porque utiliza una aproximación cuadrática y no solo lineal.
\end{respuesta}

%%%%%%%%%%%%%%%%%%%%%%%%%%%%%%%%%%%%%%%%
% Pregunta 10
%%%%%%%%%%%%%%%%%%%%%%%%%%%%%%%%%%%%%%%%

\item
Pregunta conceptual de métodos numéricos: \puntaje{3}

Explique por qué el método de Gauss-Seidel puede converger más rápido que el método de Jacobi en algunos sistemas lineales.

\[
\underline{\hspace{14cm}}
\]

\[
\underline{\hspace{14cm}}
\]

\[
\underline{\hspace{14cm}}
\]

\begin{respuesta}
El método de Jacobi calcula todos los valores nuevos usando únicamente los valores de la iteración anterior.

En cambio, el método de Gauss-Seidel utiliza de inmediato los valores que ya fueron actualizados dentro de la misma iteración. Esto permite que la información nueva se incorpore más rápido al proceso de cálculo.

Por esta razón, en algunos sistemas lineales, Gauss-Seidel puede acercarse a la solución en menos iteraciones que Jacobi.
\end{respuesta}

\end{preguntas}

\end{document}

